% concatenated from StabilityLandauMHD.tex
\documentclass[10pt]{amsart}
\usepackage[T1]{fontenc}
\usepackage[utf8]{inputenc}
\usepackage[english]{babel}
\usepackage{csquotes} 
\usepackage[a4paper,margin=3cm]{geometry}
\sloppy 
\allowdisplaybreaks
\usepackage[draft]{graphicx}
\usepackage[backend=biber,style=alphabetic,bibencoding=utf8,language=auto,autolang=other,giveninits=true,doi=false,isbn=false,url=false,maxnames=10,sorting=nty]{biblatex}
\DeclareNolabel{
  \nolabel{\regexp{[\p{P}\p{S}\p{C}\p{Z}]+}}
}
\usepackage{caption} % Better handling of empty figure captions.
\usepackage[hyperfootnotes=false]{hyperref} %See: https://github.com/plk/biblatex/issues/744
\hypersetup{
	colorlinks = true,
	linkcolor = {blue},
	urlcolor = {red},
	citecolor = {blue}
}
\usepackage[nameinlink,capitalise,sort]{cleveref}
\crefname{equation}{}{} % Makes cref behave like eqref.
%\crefname{enumi}{}{} % Avoid showing ``item (1)''
\newcommand{\crefrangeconjunction}{--}
\crefname{figure}{Figure}{Figures}


\usepackage{amsmath}
\usepackage{amsthm}
\usepackage{amssymb}
\usepackage{esint} % For the average integral
\usepackage{mathrsfs} % Just for cool sigma-algebra
\usepackage{faktor} % For quotients
\usepackage{dsfont} % For identity.
\usepackage[dvipsnames]{xcolor}
\usepackage{array}
\usepackage{hhline}
\usepackage[normalem]{ulem}
\usepackage{comment} % For comment environment
\usepackage[toc,page]{appendix} % Appendices
\usepackage{imakeidx} % Analytical index
\usepackage{xparse} % In order to use ExplSyntax (in older latex versions)
\usepackage{mathtools}
\usepackage{fancyhdr} % Header and footer of pages
\usepackage{ifthen} % Ifthen construct
\usepackage{forloop} % forloop construct
\usepackage{xstring}
\usepackage{tikz}
\usetikzlibrary{positioning, arrows.meta}
%\usepackage{pgfplots}
\usetikzlibrary{babel} % Fix buggy interaction with babel package
\usetikzlibrary{cd} % Commutative diagrams
\usepackage{tikz-cd} 
\usepackage{emptypage} % Empty page are shown completely white.
\usepackage{listings} % Source code.
\usepackage{mathabx} % for widecheck
%%\usepackage[inline]{showlabels} % To show labels' names when in draft mode.
\usepackage{subcaption}



%% THEOREMS
\theoremstyle{plain}
\newtheorem{lemma}{Lemma}[section]
\newtheorem{proposition}[lemma]{Proposition}
\newtheorem{theorem}[lemma]{Theorem}
\newtheorem{corollary}[lemma]{Corollary}
\newtheorem{property}[lemma]{Property}
\newtheorem{remark}[lemma]{Remark}

\theoremstyle{definition}
\newtheorem{definition}[lemma]{Definition}
\newtheorem{conjecture}[lemma]{Conjecture}
\newtheorem{claim}[lemma]{Claim}
\newtheorem{assumption}[lemma]{Assumption}

\theoremstyle{remark}
\newtheorem{example}[lemma]{Example}
\newtheorem{notation}[lemma]{Notation}

\numberwithin{equation}{section}


%% ------ SHORTCUTS --------
\newcommand{\BV}{\mathrm{BV}}
\newcommand{\BD}{\mathrm{BD}}
\newcommand{\bmo}{\mathrm{BMO}}
\newcommand{\vmo}{\mathrm{VMO}}
\newcommand{\loc}{\mathrm{loc}}
\newcommand{\lip}{\mathrm{Lip}}
\newcommand{\R}{\mathbb{R}}
\newcommand{\N}{\mathbb{N}}
\newcommand{\T}{\mathbb{T}}
\newcommand{\eps}{{\varepsilon}}
\newcommand{\TV}{\mathrm{TV}}
\DeclareMathOperator{\diam}{diam}
\DeclareMathOperator{\argmin}{argmin}
\DeclareMathOperator{\sign}{sign}
\DeclareMathOperator{\supp}{supp}
\DeclareMathOperator{\dive}{div}
\DeclareMathOperator{\curl}{curl}
\DeclareMathOperator*{\essinf}{ess\,\inf}
\DeclareMathOperator*{\esslim}{ess\,\lim}
\DeclareMathOperator*{\esssup}{ess\,\sup}
\newcommand{\sgn}[1]{\mathrm{sign}\left(#1\right)}
\newcommand{\wstarlim}{\overset{\ast}\rightharpoonup}
\newcommand{\wlim}{\rightharpoonup}
\newcommand{\weaks}{\stackrel{*}{\rightharpoonup}}

\newcommand{\p}{\partial}
\newcommand{\pt}{\partial_t}
\newcommand{\ptt}{\partial_{tt}^2}
\newcommand{\px}{\partial_x }

\newcommand{\bx}{\boldsymbol{x}}
\newcommand{\bb}{{\boldsymbol{b}}}

\newcommand{\be}{{\boldsymbol{e}}}

\newcommand{\bu}{{\boldsymbol{u}}}
\newcommand{\bv}{{\boldsymbol{v}}}
\newcommand{\bU}{{\boldsymbol{U}}}
\newcommand{\bw}{{\boldsymbol{w}}}
\newcommand{\bB}{{\boldsymbol{B}}}
\newcommand{\Bf}{{\boldsymbol{f}}}
\newcommand{\BF}{{\boldsymbol{F}}}

\newcommand{\vp}{\varphi}
\newcommand{\te}{\theta}
\newcommand{\ve}{\varepsilon}
% highlight helpers
\usepackage{soul}
\usepackage{xcolor}

\makeatletter
% Make soul accept common commands inside \hl (extend as needed)
\soulregister\ref7
\soulregister\eqref7
\soulregister\cref7
\soulregister\Cref7
\soulregister\autoref7

\soulregister\today7
\soulregister\cite7
\soulregister\citep7
\soulregister\citet7
\soulregister\citealp7
\soulregister\citeauthor7
\soulregister\citeyear7

% Optional: common text formatting commands used inside highlights
\soulregister\emph7
\soulregister\textbf7
\soulregister\textit7
\soulregister\texttt7
\makeatother

% Math-safe highlight used internally by \hlg
\makeatletter
\newcommand{\hlg@math}[1]{%
  \begingroup
  \setlength{\fboxsep}{0.6pt}% padding inside the highlight (tweak)
  \mathchoice
    {\colorbox{yellow!25}{$\displaystyle #1$}}%
    {\colorbox{yellow!25}{$\textstyle #1$}}%
    {\colorbox{yellow!25}{$\scriptstyle #1$}}%
    {\colorbox{yellow!25}{$\scriptscriptstyle #1$}}%
  \endgroup
}
\makeatother

\DeclareRobustCommand{\hlg}[1]{%
  \begingroup
  \sethlcolor{yellow!25}%
  \ifmmode
    \hlg@math{#1}%
  \else
    \hl{#1}%
  \fi
  \endgroup
}



\usepackage[shortlabels]{enumitem}
\newlist{thmenum}{enumerate}{1} % to be used only inside 'theorem' environments
\setlist[thmenum]{label=\textup{(\roman*)},
                  ref=\thetheorem-{(\roman*)}}
\crefname{thmenumi}{Theorem}{Theorem}

\newlist{lemenum}{enumerate}{1} % to be used only inside 'lemma' environments
\setlist[lemenum]{label={(\roman*)},
                  ref=\thelemma-{(\roman*)}}
\crefname{lemenumi}{Lemma}{Lemmas}

\newlist{propenum}{enumerate}{1} % to be used only inside 'proposition' environments
\setlist[propenum]{label={(\roman*)},
                  ref=\thelemma-{(\roman*)}}
\crefname{propenumi}{Proposition}{Propositions}

%%%% DOCUMENT %%%%
\addbibresource{LandauMHD-refs.bib}
%------------------------------------------------

\begin{document}

%%%% DOCUMENT %%%%


\title[Landau solutions for MHD]{Asymptotic stability of Landau solutions to the MHD system and energy decay}

\author[N.~De Nitti]{Nicola De Nitti}
\address[N.~De Nitti]{Università di Pisa, Dipartimento di Matematica, Largo Bruno Pontecorvo 5, 56127 Pisa, Italy.}
\email[]{nicola.denitti@unipi.it}

\author[Y.~Wang]{Yun Wang}
\address[Y.~Wang]{Soochow University, School of Mathematical Sciences, Center for dynamical systems and differential equations, 215006 Suzhou, P.\,R.~China.}
\email[]{ywang3@suda.edu.cn}

\author[S.~Zhang]{Shaoheng Zhang}
\address[S.~Zhang]{Soochow University, School of Mathematical Sciences, 215006 Suzhou, P.\,R.~China.}
\email[]{20234007008@stu.suda.edu.cn}

\keywords{Magnetohydrodynamics; asymptotic stability; Landau solutions; decay rates.}

\subjclass[2020]{%
35Q35, % PDEs in connection with fluid mechanics
35B40, % Asymptotic behavior of solutions
76W05, % Magnetohydrodynamics and electrohydrodynamics
76D05. % Navier-Stokes equations for incompressible viscous fluids
}

\begin{abstract}
%\hlg{Draft of \today.}
We consider the three-dimensional incompressible MHD system. Any weak solution satisfying a strong energy inequality is $L^2$-asymptotically stable around a Landau solution.
Under an additional integrability assumption on the initial perturbation, we also obtain an explicit algebraic decay rate for the $L^2$-norm of the velocity and magnetic perturbations.

\end{abstract}

\maketitle



\section{Introduction}
\label{sec:intro}


\subsection{The MHD and Navier--Stokes systems}\label{ssec:mhd}

The \textit{magnetohydrodynamic} (MHD) \textit{equations} describe the behavior of an electrically conducting, incompressible, viscous fluid and play a fundamental role in astrophysics, geophysics, plasma physics, and in certain industrial applications. An introduction to the topic can be found in \cite{Davidson_Book_2nd_2017}.

We consider the initial value problem for the MHD system in three space dimensions: 
\begin{align}\label{eq:IVPMHD}
\begin{cases}
\pt\bu + \bu \cdot \nabla \bu + \nabla p = \nu\, \Delta \bu + \bB \cdot \nabla \bB+\Bf, &  (\bx,t) \in \mathbb{R}^3\times(0,\infty), \\
\pt\bB + \bu \cdot \nabla \bB =  \bB \cdot \nabla \bu + \mu\, \Delta \bB, &  (\bx,t) \in \mathbb{R}^3\times(0,\infty), \\
\nabla \cdot \bu = 0, &  (\bx,t) \in \mathbb{R}^3\times(0,\infty), \\
\nabla \cdot \bB = 0, &  (\bx,t) \in \mathbb{R}^3\times(0,\infty), \\
\bu(\bx,0)=\bu_0(\bx), &  \bx \in \mathbb{R}^3,\\
\bB(\bx,0)=\bB_0(\bx), &  \bx \in \mathbb{R}^3.
\end{cases}
\end{align}
Here, the unknown functions $\bu,\bB,p$ denote the velocity, magnetic field, and pressure, respectively; the data $\bu_0,\bB_0$ and $\Bf$ represent the given initial velocity, initial magnetic field, and external force. 

The equations \cref{eq:IVPMHD}$_{1,3}$ constitute the Navier--Stokes system with an external Lorentz force, while \cref{eq:IVPMHD}$_{2,4}$ are derived from Maxwell's equations for the magnetic field $\boldsymbol{B}$. In particular, in the absence of the magnetic field (i.\,e., when $\boldsymbol{B} \equiv \boldsymbol{0}$), the MHD system \cref{eq:IVPMHD} reduces to the \emph{Navier--Stokes equations}, that is 
\begin{align}\label{eq:IVPNSE}
\begin{cases}
\pt\bu + \bu \cdot \nabla \bu + \nabla p =\nu \Delta \bu +\Bf, &  (\bx,t) \in \mathbb{R}^3\times(0,\infty), \\
\nabla \cdot \bu = 0, &  (\bx,t) \in \mathbb{R}^3\times(0,\infty), \\
\bu(\bx,0)=\bu_0(\bx), &  \bx \in \mathbb{R}^3.
\end{cases}
\end{align}

The existence of weak solutions to the Navier--Stokes system with zero external force ($\Bf\equiv \mathbf0$) goes back to Leray \cite{Leray1934}, for any divergence-free initial data $\bu_0\in L^2(\mathbb{R}^3)$; an analogous result for nonzero $\Bf$ can be found in  \cite[Chapter III]{Temam1979}. Leray also raised the question whether the $L^2$-norm of a weak solution tends to zero as $t\to\infty$. This was answered positively by Kato \cite{Kato1984} for strong solutions, and by Masuda \cite{Masuda1984} for weak solutions satisfying a strong energy inequality. Later, Schonbek \cite{Schonbek1985, Schonbek1986} established explicit decay rates for the $L^2$-norm using the Fourier splitting technique. 

For the MHD equations \cref{eq:IVPMHD}, several existence results have been established. In two spatial dimensions, when $\nu > 0$ and $\mu > 0$, global well-posedness for the Cauchy
problem has been obtained in  \cite{MR716200,MR346289} for bounded domains and in \cite{MR1616418} for the whole space. Using the  vorticity-current formulation, local-in-time existence and regularity was more recently obtained in \cite{zbMATH08005764} for $L^1$ initial data. In three dimensions, local-in-time well-posedness is contained as well in \cite{MR716200,MR346289,MR1616418} , while global-in-time well-posedness
for small initial data was later obtained in \cite{MR3780143}.\footnote{~For the inviscid MHD equations (i.\,e., $\nu = 0$ and $\mu > 0$), global well-posedness was
proved in \cite{MR1007099} in the two-dimensional setting. In three dimensions, local
well-posedness was established in \cite{MR2507450}, while global existence together with
local-in-time well-posedness in critical Besov spaces was obtained in \cite{MR4496608}.

In the case of vanishing magnetic resistivity (i.\,e., $\mu = 0$ and $\nu > 0$), local
existence and uniqueness in two-dimensional spatial domains were proved in
\cite{MR2265204,MR3217057,MR2772455}. Local existence of solutions in Besov spaces, in both
two and three dimensions, can be found in \cite{MR3415680,MR3590662}. Global well-posedness
in this regime has been achieved in \cite{MR3766969} for a special geometry, namely an
infinite slab, and for particular initial data given by a uniform non-parallel magnetic
field.

Finally, for the ideal MHD equations (i.\,e., $\mu = \nu = 0$), the local existence of strong
solutions was established in \cite{MR952901,MR1257135}; additional existence results can be found in \cite{zbMATH07327027,zbMATH05675031,zbMATH07635066,zbMATH08053923,zbMATH06383463}.}

The time-asymptotic behavior of solutions to the MHD equations has also been widely studied. In \cite{SSS1996}, lower and upper bounds are obtained for the decay rates of both the total and magnetic energies for solutions in $\mathbb{R}^n$, $2 \le n \le 4$. It is shown that weak solutions with large initial data outside the class of functions with radially equi-distributed energy exhibit algebraic decay, whereas radially equi-distributed initial data lead to exponential decay. In \cite{AgapitoSchonbek2007}, the long-time behavior of solutions in  $\mathbb{R}^n$, $2 \le n \le 3$, is investigated. In the absence of magnetic diffusion, it is proved that, provided strong bounded solutions exist, the total energy cannot exhibit asymptotic oscillations. When magnetic diffusion is present and the initial data belong only to $L^2$, solutions are shown to decay to zero without a uniform rate, and this lack of uniformity is optimal. Some decay results in exterior domains are available in \cite{MR894160}.


\subsection{Landau solutions for stationary Navier--Stokes and MHD equations}

In this work, we focus on the long-time behavior of weak solutions to the three-dimensional incompressible MHD system \cref{eq:IVPMHD} near a stationary \textit{Landau solution}.

Stationary solutions to the Navier--Stokes equations that are homogeneous of degree $-1$ in $\mathbb{R}^3\setminus\{\mathbf0\}$ are the Landau solutions introduced in \cite{Landau1944}; they are axisymmetric without swirl (see \cref{ssec:existence} for the definition and explicit formulas, or \cite[Section 8.2]{Tsai2018}).


Tian and Xin \cite{TianXin1998}, and Cannone and Karch \cite{CannoneKarch2004} proved that all $(-1)$-homogeneous, axisymmetric solutions in $C^2\left(\mathbb{R}^3 \backslash\{\mathbf0\}\right)$ are Landau solutions.  Later, in \cite{Sverak2011}, \v{S}ver\'{a}k proved the assertion without the assumption of axisymmetry. 

We stress that we focus on the three-dimensional case. Indeed, for $n = 2$, the problem of finding  $(-1)$-homogeneous solutions reduces to studying an ODE on the circle $\mathbb S^1$, where a reasonably complete description of solution can be obtained in terms of elliptic functions; under the additional constraint that $\nabla \cdot \boldsymbol{u} = 0$ across the origin, there exists (modulo rotations) a countable family of $(-1)$-homogeneous solutions that are smooth away from the origin. On the other hand, for $n \ge 4$, there are no non-trivial $(-1)$-homogeneous solutions that are smooth away from the origin. We refer to \cite{Sverak2011} and the references therein for further details.

Landau solutions plays a fundamental role in stationary Navier--Stokes flows: Korolev and \v{S}ver\'{a}k \cite{KorolevSverak2011} showed that the asymptotic leading term of small solutions in exterior domains of $\mathbb{R}^3$ must be a Landau solution; Miura and Tsai \cite{MiuraTsai2012} proved that the leading term of a point singularity like $|\bx|^{-1}$ at $\bx=\mathbf0$ is also given by a Landau solution when it is sufficiently small.

For the stationary MHD system, the pair $(\bu, \bB)$ is a $(-1)$-homogeneous, axisymmetric solution whenever $\bu$ is a Landau solution and $\bB \equiv \boldsymbol 0$. The converse direction has been investigated under different structural assumptions. For the system \textit{without magnetic diffusion} (i.\,e., $\mu=0$), Zhang \cite{Zhang2025MHDwithoutMD} showed that the converse holds provided the magnetic field $\bB$ has only the swirl component. In the case \textit{with magnetic diffusion} (i.\,e., $\mu >0$), Zhang, Wang, and Wang \cite{ZhangWangWang2026} proved that if $(\bu, \bB)$ is a smooth axisymmetric self-similar solution with $\bB = B^\theta(\rho, \varphi) \, \boldsymbol{e}_\theta$ and satisfies the pointwise decay condition $|\bu(\bx)| \leq \varepsilon |\bx|^{-1}$ for some small $\varepsilon > 0$, then $\boldsymbol{u}$ must be a Landau solution and $\boldsymbol{B}$ vanishes. 
This condition was later refined by Zhang \cite{Zhang2025MHDwithMD}, who showed that the conclusion holds under the weaker assumption that the cylindrical radial component $u^r \le \beta$ on the unit sphere for some constant $\beta < \frac{2\sqrt{2}}{3}$.


\subsection{Perturbations of the Landau solution}\label{ssec:perturb-setting}

For simplicity, let us fix $\nu=\mu=1$. Consider initial data of the form $\bu_0=\bU^\bb+\bw_0,\ \bB_0$ and $\Bf=\bb\delta_{\mathbf{0}}$, where $\bU^\bb$ is the Landau solution corresponding to a vector field $\bb$ (as defined in \cref{ssec:existence}) and $\bw_0, \bB_0 \in L^2_{\sigma}(\mathbb{R}^3)$ (cf.~\cref{ssec:weakLpspace} for the notation), provided $|\bb|>0$ is sufficiently small.
Throughout this paper, $\delta_{\mathbf{0}}$ denotes the Dirac distribution, while $
\delta, \tilde \delta,  \delta_1,\delta_2,\dots$ are positive constants.


%We establish the existence of weak solutions to problem \cref{eq:IVPMHD} for the above class of initial data and show their asymptotic convergence to the stationary solution $(\bU^\bb(\bx),\mathbf{0})$ as $t \to \infty$, yielding $L^2$-asymptotic stability.  

Denote by $(\bu(\bx,t), \bB(\bx,t), p(\bx,t))$ the solution to the MHD equations \cref{eq:IVPMHD}. Setting $\bu(\bx,t)=\bU^{\bb}(\bx)+\bw(\bx,t)$ and $p(\bx,t)=P^\bb(\bx)+\pi(\bx,t)$, where $(\bU^\bb,P^\bb)$ is the Landau solution defined in \cref{ssec:existence},  we obtain the perturbed system
\begin{align}\label{eqn:PerturbedEqn}
\begin{cases}
\pt \bw + \bw \cdot \nabla \bw + \bU^\bb\cdot \nabla\bw+\bw\cdot \nabla \bU^\bb+\nabla \pi = \Delta\bw + \bB \cdot \nabla \bB, & (\bx,t)\in \mathbb{R}^3\times(0,\infty),\\
\pt \bB+\bw \cdot \nabla \bB+ \bU^\bb \cdot \nabla \bB=\Delta \bB + \bB \cdot \nabla \bw + \bB\cdot \nabla \bU^\bb, & (\bx,t)\in \mathbb{R}^3\times(0,\infty),\\
\nabla \cdot \bw=0, & (\bx,t)\in \mathbb{R}^3\times(0,\infty),\\
\nabla \cdot \bB=0, & (\bx,t)\in \mathbb{R}^3\times(0,\infty),\\
\bw(\bx,0)=\bw_0(\bx), & \bx \in \mathbb{R}^3,\\
\bB(\bx,0)=\bB_0(\bx), & \bx \in \mathbb{R}^3.
\end{cases}
\end{align}

When $\bB \equiv \mathbf{0}$, system \cref{eqn:PerturbedEqn} reduces to the perturbed Navier--Stokes equations around the Landau solution:
\begin{align}\label{eqn:PerturbedEqnNS}
\begin{cases}
\partial_t \bw + \bw \cdot \nabla \bw + \bU^\bb \cdot \nabla \bw + \bw \cdot \nabla \bU^\bb + \nabla \pi = \Delta \bw, & (\bx,t) \in \mathbb{R}^3 \times (0,\infty), \\
\nabla \cdot \bw = 0, & (\bx,t) \in \mathbb{R}^3 \times (0,\infty), \\
\bw(\bx,0) = \bw_0(\bx), & \bx \in \mathbb{R}^3.
\end{cases}
\end{align} 
In \cite{KarchPilarczyk2011}, Karch and Pilarczyk showed that if $|\bb|$ is sufficiently small and $\bw_0 \in L^2_\sigma(\mathbb{R}^3)$, there exists at least one weak solution $\bw$ to \eqref{eqn:PerturbedEqnNS} satisfying a strong energy inequality, and that any such solution is $L^2$--asymptotically stable, i.\,e., $\|\bw(t)\|_2 \to 0$ as $t \to \infty$. Under the additional assumption $\bw_0 \in L^q(\mathbb{R}^3)$ for some $\frac{6}{5} < q < 2$, the solution also obeys the decay estimate $\|\bw(t)\|_2 \leq C t^{-\frac{3}{2}(\frac{1}{q} - \frac{1}{2})}$. 
Subsequently, Karch, Pilarczyk, and Schonbek \cite{KarchPilarczykSchonbek17} generalized the $L^2$-stability of the Landau solution to a broader class of global-in-time solutions, including singular and time-dependent flows, provided the background solution is small enough.
The energy decay rate established in \cite{KarchPilarczyk2011} was further extended by Hishida and Schonbek \cite{HishidaSchonbek16}, who proved that for any $q\in [1,2)$, there exists a constant $\delta=\delta(q)>0$ such that if $|\bb|\le \delta$, the typical decay rate $t^{-\frac{3}{2}(\frac{1}{q}-\frac{1}{2})}$ also holds, thus extending the admissible range of $q$. 

\v{S}ver\'{a}k \cite{Sverak2011} proved that every $(-1)$-homogeneous non-zero solution of the stationary Navier--Stokes equations in $C^2(\mathbb{R}^3\setminus\{\mathbf0\})$ is a Landau solution.  
In \cite{LLY18I,LLY18II,LLY19III}, the authors studied $(-1)$-homogeneous axisymmetric solutions that are smooth away from the symmetry axis $\{x_1 = x_2 = 0\}$ and may be singular along the entire axis. In this class, all no-swirl solutions were classified in \cite{LLY18I,LLY18II}, while solutions with non-zero swirl were constructed in \cite{LLY18I,LLY19III}.  Among these solutions, the least singular ones (those with sufficiently mild singularities along the axis) were shown in \cite{LiYan2021} to be asymptotically stable under $L^2$-perturbations, thereby extending the stability result of \cite{KarchPilarczyk2011} beyond the Landau case.  

In \cite{LZZ23}, the $L^3$-asymptotic stability of Landau solutions was studied:  if $\boldsymbol{w}_0$ and the background Landau solution are sufficiently small, then there exists a unique global solution with $\|\boldsymbol{w}(t)\|_{3} \to 0$ as $t \to \infty$.   On the other hand, in \cite{zbMATH08070132}, the authors showed that the asymptotic stability around Landau solutions may fail for some initial perturbations in $L^{3,\infty}$ (regardless of how small the Landau solution or the initial perturbations are in $L^{3,\infty}$). 

Finally, in \cite{zbMATH07471753}, the authors studied the stability and asymptotic behavior of singular solutions to the 3D incompressible Navier--Stokes equations with singular external forces which are either singular finite measures or more general tempered distributions
with bounded Fourier transforms (including, in particular, the Landau solution).

We now define the notion of a weak solution to \cref{eqn:PerturbedEqn}.
For initial data $\bw_0, \bB_0\in L^2_{\sigma}$, $(\bw,\bB)=(\bw(\bx,t),\bB(\bx,t))$ is called a \textit{weak solution} if $\bw,\bB\in X$, where
\[
X\coloneqq C_w \big( [0,\infty); L^2_\sigma (\R^3)\big)\cap L^2\big([0,\infty); \dot H_\sigma^1(\R^3)\big),
\]
and if for all $t\geq s\geq 0$ and all test functions $\boldsymbol{\vp}$ belonging to
\begin{align}\label{eqn:testfcnspace}
\boldsymbol{\vp} \in C\big( [0, \infty); L^2_\sigma \cap L^3 \big), \quad 
\nabla \boldsymbol{\vp} \in L^2_{\mathrm{loc}}\big( [0,\infty); L^2 \big), \quad 
\partial_\tau \boldsymbol{\vp} \in L^2_{\mathrm{loc}}\big( [0,\infty); (H^1_\sigma)^* \big),
\end{align}
the following identities hold:
\begin{equation}\label{eqn:weak}
\begin{aligned}
&\langle \bw(t), \boldsymbol{\varphi}(t)  \rangle
+\int_s^t \Big[ 
\langle \bw\cdot  \nabla \bw, \boldsymbol{\varphi} \rangle
- \langle \bB\cdot  \nabla \bB, \boldsymbol{\varphi} \rangle
+ a_1(\bw,\boldsymbol{\vp})
\Big]\, \mathrm d\tau 
\\ &\qquad= \langle \bw(s), \boldsymbol{\varphi}(s)  \rangle 
+ \int_s^t \langle \bw, \boldsymbol{\varphi}_\tau \rangle \, \mathrm d\tau,\\
&\langle \bB(t), \boldsymbol{\varphi}(t)  \rangle
+\int_s^t \Big[ 
\langle \bw\cdot  \nabla \bB, \boldsymbol{\varphi} \rangle
-\langle \bB\cdot  \nabla \bw, \boldsymbol{\varphi} \rangle
+ a_2(\bB,\boldsymbol{\vp})
 \Big]\, \mathrm d\tau 
 \\ &\qquad= \langle \bB(s), \boldsymbol{\varphi}(s)  \rangle 
 + \int_s^t \langle \bB, \boldsymbol{\varphi}_\tau \rangle \, \mathrm d\tau.
\end{aligned}
\end{equation}
Here, $\langle \cdot,\cdot \rangle$ denotes the standard $L^2$-inner product, and $a_1(\cdot,\cdot)$ and $a_2(\cdot,\cdot)$ are defined in \cref{eqn:bilinearform}.
The test function space is a special case of the more general framework introduced in \cite[eq.~(2.16)]{HishidaSchonbek16}. 


\subsection{Main results}\label{ssec:main}
This work investigates the long-time behavior and decay rates of weak solutions to \cref{eqn:PerturbedEqn} that obey a strong energy inequality of the form
\begin{align}\label{eqn:energyineq1}
\begin{aligned}
&\|\bw(t)\|_2^2 + \|\bB(t)\|_2^2 + 2 \int_s^t \bigl( \| \nabla \bw(\tau)\|_2^2 + \| \nabla \bB(\tau)\|_2^2 \bigr) \,\mathrm d \tau\\
& \le \|\bw(s)\|_2^2 + \|\bB(s)\|_2^2 - 2 \int_s^t \bigl( \langle (\bw \cdot \nabla) \bU^\bb, \bw \rangle - \langle (\bB \cdot \nabla) \bU^\bb, \bB \rangle \bigr) \,\mathrm d \tau,
\end{aligned}
\end{align}
which holds for almost every $s \ge 0$ (including $s=0$) and all $t \ge s$. We note that, as in \cite[Theorem 2.1]{KarchPilarczyk2011}, Galerkin's method yields the existence of (at least) one weak solution to \cref{eqn:PerturbedEqn} satisfying \cref{eqn:energyineq1} provided $\bU^\bb$ is small.

Our first main result is as follows.
\begin{theorem}[Asymptotic stability]\label{thm:stability}
There exists a constant $\delta>0$ such that if $0<|\bb|\le \delta$, then every weak solution $(\bw(t),\bB(t))$ to \cref{eqn:PerturbedEqn} with \cref{eqn:energyineq1} satisfies
\begin{align*}
    \lim\limits_{t\to \infty} \|\bw(t)\|_2=\lim\limits_{t\to \infty} \|\bB(t)\|_2=0,
\end{align*}
provided $\bw_0, \bB_0 \in L^2_{\sigma}(\R^3)$.
\end{theorem}

Under additional integrability assumptions on initial data, we also get the decay rate of $\|\bw(t)\|_2$ and $\|\bB(t)\|_2$.

\begin{theorem}[Energy decay rate]\label{thm:decay}
    Let $1\le q<2$. There exists a constant $\tilde\delta=\tilde\delta(q)>0$ such that if $0<|\bb|\le \tilde\delta$, then every weak solution $(\bw(t),\bB(t))$ to \cref{eqn:PerturbedEqn} with \cref{eqn:energyineq1}  satisfies
\begin{align*}
    \|\bw(t)\|_2 \le C(t+1)^{-\frac{3}{2}(\frac{1}{q}-\frac{1}{2})},\quad 
    \|\bB(t)\|_2 \le C(t+1)^{-\frac{3}{2}(\frac{1}{q}-\frac{1}{2})},
\end{align*}
provided $\bw_0, \bB_0 \in L^q(\R^3)\cap L^2_{\sigma}(\R^3)$.
\end{theorem}

These results extend the analysis of decay properties for weak solutions of \cref{eqn:PerturbedEqnNS} satisfying a strong energy inequality analogous to \cref{eqn:energyineq1}, previously carried out in \cite{HishidaSchonbek16}. Unlike \cite{HishidaSchonbek16}, we do not use the Fourier splitting method to obtain the energy decay; instead, we present a streamlined argument that combines ideas from \cite{HishidaSchonbek16}, \cite{KarchPilarczyk2011}, and \cite{BorchersMiyakawa92}. In particular, inspired by \cite{KarchPilarczyk2011}, we start from estimates of the type 
\begin{align*}
\|\bw(s)\|_2 
&\leq \|e^{-sL_1} \bw_0\|_2 
+ \int_0^s \bigl\|e^{-(s-\tau)L_1} \mathbb{P} \nabla\cdot (\bw \otimes \bw - \bB \otimes \bB)\bigr\|_2 \, \mathrm{d}\tau, \\
\|\bB(s)\|_2 
&\leq \|e^{-sL_2} \bB_0\|_2 
+ \int_0^s \bigl\|e^{-(s-\tau)L_2} \mathbb{P} \nabla\cdot (\bw \otimes \bB - \bB \otimes \bw)\bigr\|_2 \, \mathrm{d}\tau,
\end{align*}
(cf.~\crefrange{eqn:estimate2-1}{eqn:estimate2-2} below) and use them to derive the energy decay of solutions. The main difference from \cite{KarchPilarczyk2011} lies in our treatment of the linear and nonlinear terms: we employ the more flexible $L^q$--$L^2$ estimates from \cite{HishidaSchonbek16} (see, ~\crefrange{eqn:decay1}{eqn:decay2}), which yield a more direct proof.

\subsection{Outline} 
The paper is organized as follows. \cref{sec:prelim} collects some preliminary lemmas that will be used. 
In particular, in \cref{ssec:existence}, we recall the explicit form of the Landau solutions to the stationary Navier--Stokes equations and state Hardy's inequality; in \cref{ssec:weakLpspace}, we introduce the weak-$L^p$ (Marcinkiewicz) spaces and present the weak versions of Young's and H\"older's inequalities.
In \cref{sec:linear}, we study the linearized perturbed system, proving that the linear operators generate analytic semigroups and stating the decay rate of the evolution operator. Finally, the proofs of \crefrange{thm:stability}{thm:decay} are carried out in \cref{sec:proofs}.


\section{Preliminaries}
\label{sec:prelim}

\subsection{Laudau solutions and Hardy-type inequality}
\label{ssec:existence}

This subsection is devoted to recalling the Landau solutions to the stationary Navier--Stokes equations and Hardy's inequality.

We begin with some basic facts about the Landau solutions (see e.\,g., \cite[Section 8.2]{Tsai2018}). For each nonzero $\bb \in \mathbb{R}^3$, there exists a unique Landau solution $\bU^\bb$ and an associated pressure $P^\bb$, homogeneous of degrees $-1$ and $-2$ respectively, that solve the stationary Navier--Stokes equations
\begin{align}\label{eq:NS}
\begin{cases}
\bU^\bb \cdot \nabla \bU^\bb + \nabla P^\bb = \Delta \bU^\bb, & \bx \in \R^3\setminus \{\mathbf0\}, \\
\nabla \cdot \bU^\bb = 0, & \bx \in \R^3\setminus \{\mathbf0\},
\end{cases}
\end{align}
and are smooth in $\mathbb{R}^3 \setminus \{\mathbf0\}$. Moreover, in the sense of distributions in $\mathbb{R}^3$, they satisfy
\begin{align*}
-\Delta \bU^\bb + \nabla \cdot (\bU^\bb \otimes \bU^\bb) + \nabla P^\bb = \bb \, \delta_{\mathbf{0}}, \qquad 
\nabla \cdot  \bU^\bb = 0,
\end{align*}
where $\delta_{\mathbf{0}}$ denotes the Dirac distribution.

When $\bb = (0,0,\beta)$ with $\beta>0$, the Landau solution admits an explicit representation in spherical coordinates. Let $\bx = (\rho\sin\varphi\cos\theta,\rho\sin\varphi\sin\theta, \rho\cos\varphi)$ and define the orthonormal basis vectors
\[
\be_{\rho}=\frac{\bx}{\rho}, \quad \be_{\theta}=(-\sin\theta, \cos\theta,0), \quad \be_{\varphi}=\be_{\theta}\times \be_{\rho}.
\]
Then for any $A>1$,
\begin{align}\label{eqn:Landausolution}
    \bU^\bb =\frac{2}{\rho} 
    \left[
    \left(\frac{A^2-1}{(A-\cos\varphi)^2}-1\right)\be_{\rho}
    +\frac{-\sin\varphi}{A-\cos\varphi}\be_{\varphi}
    \right], \quad
    P^\bb =\frac{4(A \cos\varphi-1)}{\rho^2 (A-\cos\varphi)^2},
\end{align}
where the parameter $\beta$ is related to $A$ by
\[
\beta = 16 \pi 
\left[A + \frac 12 A^2 \log \left(\frac {A-1}{A+1}\right) + \frac{4A}{3(A^2-1)}\right].
\]
The corresponding Landau solution for a general $\bb$ is obtained by rotation.


There exist $\ve_0>0$ and $C_0>0$ such that for any $\bb \in \R^3$ with $0 < |\bb| < \ve_0$, 
\begin{align}\label{eqn:LandauEst}
    |\bU^\bb(\bx)|\le\frac{C_0|\bb|}{|\bx|},\quad |\nabla \bU^\bb(\bx)|\le\frac{C_0|\bb|}{|\bx|^2},\quad \text{in }\R^3\setminus\{\mathbf0\}.
\end{align}
Indeed, a direct computation shows that for $\bb=(0,0,\beta)$ with sufficiently small $|\beta|>0$,
%In fact, a direct computation shows that 
\[
|\bU^\bb(\bx)| \lesssim \frac{1}{|A|} \frac{1}{|\bx|},\quad
|\nabla \bU^\bb(\bx)| \lesssim \frac{1}{|A|} \frac{1}{|\bx|^2}.
\]
%when $\bb=(0,0,\beta)$ with sufficiently small $|\beta|>0$.     
Combining this with $\frac{\beta(A)}{16 \pi}=\frac{1}{A}+o(\frac{1}{A})$ as $A\to \infty$ yields \eqref{eqn:LandauEst}.

A vector field $\bv$ is called \emph{axisymmetric} if it is of the form
\[
\bv=v^\rho(\rho,\vp) \be_\rho + v^\theta(\rho,\vp) \be_\theta + v^\vp(\rho,\vp) \be_\vp,
\]
with $v^\rho, v^\theta, v^\vp$ independent of $\theta$. 
It is called \emph{without swirl} if $v^\theta = 0$.


Finally, we recall the following \textit{Hardy-type inequality}:
\begin{align}\label{eqn:Hardy}
    \int_{\R^3} \frac{|w(\bx)|^2}{|\bx|^2}\, \mathrm{d}x \le 4 \int_{\R^3} |\nabla w(\bx)|^2 \, \mathrm{d}x \quad \text{for all } w\in \dot{H}^1(\R^3),
\end{align}
a proof of which is given in \cite[Chapter I.6]{Leray1934}.\footnote{~We remark that \cref{eqn:Hardy} is valid when $n\ge 3$, with the constant $4$ replaced by $\frac{4}{(n-2)^2}$  (see for example \cite[Theorem 2.1]{AGG2006}); this restriction constitutes a technical difficulty in the analysis of the stability of $(-1)$-homogeneous solutions in $\mathbb R^2$.}

\subsection{Notations and inequalities in weak \texorpdfstring{$L^p$}{Lp} spaces}
\label{ssec:weakLpspace}

The following notation will be used throughout the paper. Let $L^p_\sigma(\mathbb{R}^3)$ denote the completion of $C_{0,\sigma}^{\infty}(\mathbb{R}^3)$ in the $L^p$-norm $\|\cdot\|_{L^p(\mathbb{R}^3)}$, where $C_{0,\sigma}^{\infty}(\mathbb{R}^3)$ is the space of smooth, compactly supported, divergence-free vector fields.

For $1<p<\infty$, let $L^{p,\infty}=L^{p,\infty}(\mathbb{R})$ denote the \textit{weak Marcinkiewicz spaces}, i.\,e., Banach spaces of measurable functions on $\mathbb{R}$ equipped with the norm
\begin{align}\label{eqn:weakLpnorm}
\|f\|_{p,\infty}= \sup
\left\{|E|^{-1+\frac{1}{p}} \int_E |f(s)| \, \mathrm{d}s: E \in \mathcal{B}\right\}.
\end{align}
Here $|E|$ denotes the Lebesgue measure of a measurable set $E$, and $\mathcal{B}$ is the collection of all Borel sets with finite positive measure.


One can verify that $|\cdot|^{-\frac{1}{p}} \in L^{p,\infty}(0,\infty)$ for $p>1$. This follows from the Hardy--Littlewood inequality $\int_E s^{-\frac{1}{p}} \, \mathrm{d}s \le \int_0^{|E|} s^{-\frac{1}{p}} \, \mathrm{d}s$ together with \cref{eqn:weakLpnorm}.
Recall also the well-known embedding $L^p \subset L^{p,\infty}$, which is a consequence of H\"older's inequality and \cref{eqn:weakLpnorm}.

The following lemma provides weak versions of Young's and H\"older's inequalities (see  \cite[Lemma 2.2]{BorchersMiyakawa92} for a detailed proof).

\begin{lemma} 
Let $f\in L^{p,\infty}$ and $g\in L^{q,\infty}$ with $1<p,q<\infty$.
{
\renewcommand{\descriptionlabel}[1]{\emph{#1}}
\begin{description}
\item[Weak Young's inequality:] If $1+\frac{1}{r}=\frac{1}{p}+\frac{1}{q}$ with $1<r<\infty$, then the convolution $f\ast g$ belongs to $L^{r,\infty}$ and satisfies
\begin{align}\label{eqn:Youngineq}
\|f \ast g\|_{r,\infty} \le C_{p,q} \|f\|_{p,\infty} \|g\|_{q,\infty},
\end{align}
where $C_{p,q}>0$ is a constant depending only on $p$ and $q$.

\item[Weak H\"older's inequality:] If $\frac{1}{r}=\frac{1}{p}+\frac{1}{q}$ with $1<r<\infty$, then the pointwise product $fg$ belongs to $L^{r,\infty}$ and satisfies
\begin{align}\label{eqn:Holderineq}
\|fg\|_{r,\infty} \le C_{p,q} \|f\|_{p,\infty} \|g\|_{q,\infty},
\end{align}
where $C_{p,q}>0$ is a constant depending only on $p$ and $q$.
\end{description}
}
\end{lemma}



\section{Linearized flow}
\label{sec:linear}



In this section, we consider the following linearized perturbed equations:
\begin{align}\label{eqn:LPE}
\begin{cases}
\pt \bw - \Delta\bw +  \bU^\bb\cdot \nabla\bw+\bw\cdot \nabla \bU^\bb+\nabla \pi = 0, 
&(\bx,t)\in \mathbb{R}^3\times(0,\infty),\\
\pt \bB - \Delta \bB + \bU^\bb \cdot \nabla \bB-\bB\cdot \nabla \bU^\bb =0,
& (\bx,t)\in \mathbb{R}^3\times(0,\infty),\\
\nabla \cdot \bw =0, 
& (\bx,t)\in \mathbb{R}^3\times(0,\infty),\\
\nabla \cdot \bB=0, & (\bx,t)\in \mathbb{R}^3\times(0,\infty),\\
\bw(\bx,0)=\bw_0(\bx), & \bx \in \mathbb{R}^3,\\
\bB(\bx,0)=\bB_0(\bx), & \bx \in \mathbb{R}^3.
\end{cases}
\end{align}
The \emph{Leray projector} $\mathbb P$ is given by $(\delta_{jk}+R_jR_k)_{1\le j,k\le n}$, where $R_j$ signifies the \emph{Riesz transform}. When $1<q<\infty$, this operator remains bounded on $L^q$ and acts as a projection onto $L^q_{\sigma}$, corresponding to the \textit{Helmholtz decomposition} $L^q = L^q_{\sigma} \oplus \{\nabla p : \ p \in L^q_{\loc}\}$ for $1<q<\infty$ (see, e.\,g., \cite[Chapter III]{Galdi11} or \cite[Section 1]{Tsai2018}). 


Applying the Leray projector $\mathbb P$ to \cref{eqn:LPE} formally yields
%Using the Leray projector $\mathbb P$, we can formally transform \cref{eqn:LPE} into 
\begin{align}\label{eqn:LPE2}
\begin{cases}
\p_t \bw + L_1 \bw =0, & t \in   (0,\infty), \\ \bw(0)=\bw_0,\\
\p_t \bB + L_2 \bB =0, & t \in (0,\infty), \\ \bB(0)=\bB_0,
\end{cases}   
\end{align}
where 
\begin{align}\label{eqn:linearoperator}
\begin{aligned}
L_1 \bw&\coloneqq \mathbb P \bigl[- \Delta\bw +  \bU^\bb\cdot \nabla\bw + \nabla\cdot  (\bw \otimes \bU^\bb)\bigr],\\
L_2 \bB&\coloneqq \mathbb P \bigl[- \Delta \bB + \bU^\bb \cdot \nabla \bB-\nabla\cdot (\bB \otimes \bU^\bb )\bigr].
\end{aligned}
\end{align}


Denote the adjoint linear operators by
\begin{align}\label{eqn:adjlinoperator}
\begin{aligned}
L_1^*\bv &=\mathbb P \bigl[-\Delta \bv-(\bU^\bb\cdot \nabla)\bv-\sum_{j=1}^3 U_j^\bb \nabla v_j\bigr],\\
L_2^*\bv &=\mathbb P \bigl[-\Delta \bv-(\bU^\bb\cdot \nabla)\bv+\sum_{j=1}^3 U_j^\bb \nabla v_j\bigr],
\end{aligned}
\end{align}
and the corresponding bilinear forms are
\begin{equation}
\begin{aligned}
\label{eqn:bilinearform}
a_1(\bw,\bv)&=\langle \nabla \bw , \nabla \bv \rangle 
    + \langle (\bU^\bb \cdot \nabla) \bw , \bv \rangle 
    +\langle (\bw \cdot \nabla) \bU^\bb, \bv \rangle,\\ 
a_2(\bB,\bv)&=\langle \nabla \bB,\nabla \bv \rangle 
    + \langle (\bU^\bb \cdot \nabla) \bB , \bv \rangle 
    -\langle (\bB \cdot \nabla) \bU^\bb, \bv \rangle.    
\end{aligned}
\end{equation}
Then
\[
a_1(\bw,\bv)=\langle L_1 \bw,\bv \rangle =\langle  \bw, L_1^*\bv \rangle,\quad 
a_2(\bB,\bv)=\langle L_2 \bB,\bv \rangle =\langle  \bB, L_2^*\bv\rangle.
\]

Let us recall a criterion due to Lions (see \cite[Proposition 1.2, Remarque 1.1]{Lions61}), which gives sufficient conditions for an operator to generate a holomorphic semigroup on a Hilbert space.

\begin{proposition}\label{prop:Lions}
Let $\mathcal{H}$ be a Hilbert space and let $\mathcal{V}\subset \mathcal{H}$ be a dense subspace. Assume that $\mathcal{V}$ is a Hilbert space with the inner product $(\cdot , \cdot )_{\mathcal{V}}$ and with the norm $\| \cdot \|_{\mathcal{V}}$ such that for a constant $C>0$ we have $ \|x\|_{\mathcal{H}}\leqslant C\|x\|_{\mathcal{V}}$ for all $x\in \mathcal{V}$. Let $a(x,y)$ be a bounded sesquilinear form on $\mathcal{V}$, which defines an operator $A: \mathcal{D}(A)\rightarrow \mathcal{H}$ as follows
\begin{equation*}
   \mathcal{D}(A)=\{ z\in \mathcal{V} :\ |a(z,v)|\leqslant C \| v\|_{\mathcal{H}}, v\in \mathcal{V}\}, \quad (Az,v)_{\mathcal{H}}=a(z,v).
\end{equation*}
Suppose that for some $\alpha >0$ and $\lambda_0 \in \mathbb R$ we have 
\begin{equation}\label{eqn:Lions}
  \alpha \| z\|_{\mathcal{V}}^2\leqslant \text{\rm Re} \ a(z,z)+\lambda_0\|z\|_{\mathcal{H}}^2.
\end{equation}
Then $-A$ is the infinitesimal generator of a strongly continuous semigroup of linear operators on $\mathcal{H}$ which is holomorphic in a sector $S_\ve\coloneqq \{s\in \mathbb C: \ |\text{\rm Arg}\ s|<\ve\}$ for some $\ve >0$.
\end{proposition}

The next theorem follows directly from \cref{prop:Lions} and the Hardy inequality \eqref{eqn:Hardy}. It establishes that for sufficiently small $|\bb|$, the operators $L_i,\,L_i^*$ $(i=1,2)$ are the infinitesimal generators of analytic semigroups of linear operators on $L^2_{\sigma}(\mathbb R^3)$.

\begin{theorem}\label{thm:analyticsemigroup}
Let $\mathbf0\neq \bb \in \R^3$ satisfy $|\bb|\le \min\{\ve_0,\frac{1}{6C_0}\}$, where $C_0$ is defined in \cref{eqn:LandauEst}. Then the operators $L_i$ and $L_i^*$ (for $i=1,2$)  defined in \cref{eqn:linearoperator} and \cref{eqn:adjlinoperator} are the infinitesimal generators of analytic semigroups of linear operators on $L^2_{\sigma}(\mathbb R^3)$ which are holomorphic in a sector $ \{s\in \mathbb C: \ |\text{\rm Arg}\ s|<\ve\}$ for some $\ve=\ve(|\bb|) >0$..
\end{theorem}

\begin{proof}
The argument follows that of \cite[Theorem 4.2]{KarchPilarczyk2011}; we outline the details for completeness.
%We prove the theorem via an argument analogous to that of \cite[Theorem 4.2]{KarchPilarczyk2011}. For completeness, we outline the details.
Take $\mathcal{H}=L^2_\sigma (\R^3)$ and $\mathcal{V}=H^1_\sigma (\R^3)$ in \cref{prop:Lions}. Then \cref{eqn:LandauEst} and Hardy's inequality \cref{eqn:Hardy} yield that
\[
|\langle (\bU^\bb \cdot \nabla) \bw , \bv \rangle| \le C_0  |\bb| \| \nabla \bw \|_2 \left\| \frac{\bv}{|\cdot|} \right\|_2
\le 2 C_0  |\bb| \| \nabla \bw \|_2 \| \nabla \bv \|_2,
\]
and 
\[
|\langle (\bw \cdot \nabla) \bU^\bb, \bv \rangle| \le C_0 |\bb| \left\| \frac{\bw}{|\cdot|} \right\|_2 
\left\| \frac{\bv}{|\cdot|} \right\|_2 
\le 4C_0 |\bb|  \| \nabla \bw \|_2 \| \nabla \bv \|_2,
\]
provided $|\bb| \le \ve_0$.
Consequently,
\begin{equation}
\begin{aligned}\label{eqn:bddcoercive}
a_1(\bw,\bv) &\le (1+6C_0|\bb|) \|\nabla \bw \|_2 \|\nabla \bv \|_2,\\ 
a_1(\bw,\bw) &\ge (1-6C_0|\bb|) \|\nabla \bw \|_2^2.
\end{aligned}   
\end{equation}
Hence \cref{eqn:Lions} is satisfied provided $|\bb| \le \min \{\ve_0,\frac{1}{6C_0}\}$, which gives the desired conclusion for $L_1$ and $L_1^*$. The same argument applies to $L_2$ and $L_2^*$.
%The conclusion for $L_2$ and $L_2^*$ follows by analogous arguments.
\end{proof}


We now state a corollary that collects standard properties of generators of analytic semigroups.
Although stated for $L_i$ ($i=1,2$), they also hold for the adjoint operators $L_i^*$.


\begin{corollary}
Under the assumptions of \cref{thm:analyticsemigroup}, the following statements hold.
\begin{enumerate}[label={(\roman*)}]
\item  For $i=1,2$, we have 
\begin{align}\label{eqn:corollary1}
(1-6C_0|\bb|)^{\frac{1}{2}} \| \nabla \bw \|_2 \le 
\| L_i^{\frac{1}{2}} \bw \|_2 \le (1+6C_0|\bb|)^{\frac{1}{2}} \| \nabla \bw \|_2  
\end{align}
for all $\bw\in \dot H^1_\sigma(\R^3)$;

\item For $i=1,2$, we have 
\begin{align}\label{eqn:corollary2}
\|e^{-tL_i}\bw_0\|_2 \le \| \bw_0\|_2  
\end{align}
for all $\bw_0\in L^2_\sigma(\R^3)$ and $t>0$.
\end{enumerate}
\end{corollary}

\begin{proof}
We prove the statements for $L_1$; the case of $L_2$ is analogous.
\begin{enumerate}[label={(\roman*)}]
\item By definition of a square root of nonnegative operators, $\|L_1^{\frac{1}{2}} \bw \|_2^2 = a_1(\bw, \bw)$. 
Hence the desired inequality follows from \cref{eqn:bddcoercive}. The same argument applies to $L_2$
\item The proof follows the same lines as for the Navier--Stokes case in \cite[Corollary 4.4]{KarchPilarczyk2011}. Multiplying equation \cref{eqn:LPE}$_1$ by $\bw$ and integrating over $\R^3$, we obtain 
\begin{equation*}
\frac{1}{2}\frac{\mathrm d}{\mathrm d t}\| \bw(t) \|^2_2+\| \nabla \bw(t)\|_2^2 
%+\langle (\bU^\bb \cdot \nabla )\bw, \bw \rangle
+\langle (\bw \cdot \nabla )\bU^\bb, \bw \rangle
=0.
\end{equation*}
Here we used $\langle (\bU^\bb \cdot \nabla )\bw, \bw \rangle=0$ because of $\operatorname{div}\, \bU^\bb=0$.
Therefore, Hardy's inequality \cref{eqn:Hardy} yields
\begin{equation*} 
   \frac{1}{2}\frac{\mathrm d}{\mathrm d t}\| \bw(t)\|^2_2 + \big( 1- 4C_0 |\bb| \big)\| \nabla \bw(t)\|_2^2\le 0,
\end{equation*}
where $1- 4C_0 |\bb|\ge 0$ because $|\bb|\le \frac{1}{6C_0}$. 
Integrating the above inequality from $0$ to $t$ gives the desired estimate.
%Now, it is sufficient to integrate from $0$ to $t$ to obtain the inequality. 
\end{enumerate}
\end{proof}

The decay estimates below for the semigroups $e^{-tL_i}$ were first established by Hishida and Schonbek in  \cite{HishidaSchonbek16}, and play a key role in the proofs of \crefrange{thm:stability}{thm:decay}. Indeed, the linearization decouples the system \cref{eqn:LPE2}, enabling us to argue similarly to  \cite{HishidaSchonbek16}. 
Their proof is based on the Duhamel formulation of the linearized equation. Treating the term $\bu \otimes \bU^\bb + \bU^\bb \otimes \bu$ as a perturbation and applying the key estimate by Yamazaki \cite[Corollary 2.3(2)]{MR1777114} in Lorentz spaces, they set up a contraction mapping in an appropriately weighted $L^p$ space, under the condition that the basic flow $\bU^\bb$ is sufficiently small in the critical space $L^{n,\infty}$. 



\begin{proposition}\label{prop:decay} 
Let $L_i$, with $i =1,\,2$, be the linear operators defined in \cref{eqn:linearoperator}. Then the following propositions hold.
\begin{propenum}
\item \label{it:lmdecay1} For any $1 \le q \le 2$, there exists a constant $\delta_1=\delta_1(q)>0$ (independent of $i$) such that if $0<|\bb|\le \delta_1$, then the semigroup $e^{-tL_i}$ satisfies
\begin{align}\label{eqn:decay1}
\|e^{-tL_i} \Bf \|_2 \le Ct^{-\frac{3}{2}(\frac{1}{q}-\frac{1}{2})} \|\Bf\|_q,\quad \text{for all } t>0, \ \Bf \in  L^q_{\sigma}(\mathbb R^3). 
\end{align}

\item \label{it:lmdecay2}
For any $\frac{6}{5} < q \le 2$, there exists $\delta_2=\delta_2(q)>0$ (independent of $i$) such that if $0<|\bb|\le \delta_2$, then the semigroup $e^{-tL_i}$ satisfies
\begin{align}\label{eqn:decay2}
\|e^{-tL_i}\mathbb P \nabla\cdot \BF\|_2 
\le Ct^{ -\frac{3}{2}(\frac{1}{q}-\frac{1}{2}) -\frac{1}{2} } \| \BF \|_q,\quad \text{for all } t>0, \  \BF\in L^q(\mathbb R^3)^{3\times 3}. 
\end{align}
\end{propenum}
\end{proposition}

\begin{proof} We prove the two claims by relying on the results in \cite{HishidaSchonbek16}. 
\begin{enumerate}[label={(\roman*)}]
    \item Since equations \eqref{eqn:LPE2}$_1$ and \eqref{eqn:linearoperator} coincide with eq.~(2.7) in \cite{HishidaSchonbek16}, the decay estimate for $e^{-tL_1}$ follows directly from \cite[Theorem 2.2]{HishidaSchonbek16} by choosing $r_0 = 3$ and $r = 2$. The proof for $e^{-tL_2}$ is analogous and will be omitted.
\item The estimate \cref{eqn:decay2} is implicitly contained in the proof of  \cite[Proposition~4.6]{HishidaSchonbek16}.
For any $1<q<\infty$, fix an exponent $r_2$ satisfying  
\[
\max\{q,\tfrac32\}<r_2<\infty,\qquad \frac1q-\frac1{r_2}<\frac13.
\] 
Such an $r_2$ exists and is determined by $q$. From the proof of Proposition~4.6 in \cite{HishidaSchonbek16}, there exists a constant $\bar \delta=\bar \delta(q)>0$ such that if $|\bb|\le \bar\delta$, then the following estimates hold for all $t>0$ and $\BF\in L^q(\mathbb R^3)^{3\times3}$:
\begin{align*}
\|e^{-tL_1}\mathbb P\operatorname{div}\BF\|_q &\le C t^{-1/2}\|\BF\|_q,\\
\|e^{-tL_1}\mathbb P\operatorname{div}\BF\|_{r_2,\infty} &\le C t^{-\frac32\bigl(\frac1q-\frac1{r_2}\bigr)-\frac12}\|\BF\|_q.
\end{align*}
% \[
% \|e^{-tL_1}\mathbb P\operatorname{div}\BF\|_q \le C t^{-1/2}\|\BF\|_q,
% \]
% \[
% \|e^{-tL_1}\mathbb P\operatorname{div}\BF\|_{r_2,\infty} \le C t^{-\frac32\bigl(\frac1q-\frac1{r_2}\bigr)-\frac12}\|\BF\|_q.
% \]
One can verify that $q<2<r_2$ provided $\frac65<q<2$.
%Now assume $\frac65<q<2$. Then we have $q<2<r_2$. 
By real interpolation of Lorentz spaces (see, e.\,g., \cite[Theorem 5.3.1]{BerghLofstrom76}),
\[
(L^{q,\infty}, L^{r_2,\infty})_{\theta,2}=L^{2,2}=L^2,
\]
where $\theta\in(0,1)$ is defined by $\frac12 = \frac{1-\theta}{q} + \frac{\theta}{r_2}$.
Consequently,
\[
\|e^{-tL_1}\mathbb P\operatorname{div}\BF\|_2 \le C \|e^{-tL_1}\mathbb P\operatorname{div}\BF\|_{q,\infty}^{1-\theta}
\|e^{-tL_1}\mathbb P\operatorname{div}\BF\|_{r_2,\infty}^{\theta}.
\]
Using the two estimates above and the elementary embedding $\|\cdot\|_{q,\infty}\le C\|\cdot\|_q$, we obtain
\[
\|e^{-tL_1}\mathbb P\operatorname{div}\BF\|_2 
%\le C t^{-\frac12(1-\theta) -\theta\bigl(\frac32(\frac1q-\frac1{r_2})+\frac12\bigr)} \|\BF\|_q
\le  C t^{-\frac32\bigl(\frac1q-\frac12\bigr)-\frac12}\|\BF\|_q.
\]
The proof for $e^{-tL_2}$ is analogous and is omitted. 
Let $\bar{\bar\delta}(q)$ denote the corresponding constant for $L_2$ and set 
$\delta_2(q)=\min\{\bar\delta(q),\bar{\bar\delta}(q)\}$. 
This completes the proof.
\end{enumerate}
\end{proof}


\section{Asymptotic stability of weak solutions}
\label{sec:proofs}

In this section, we prove \cref{thm:stability} and \cref{thm:decay}. 
 Given a weak solution $(\bw,\bB)$ to \cref{eqn:PerturbedEqn}, a duality argument yields, for all $s \ge 0$,
\begin{align*}
\|\bw(s)\|_2 &\leq \|e^{-sL_1} \bw_0\|_2 + \int_0^s \bigl\|e^{-(s-\tau)L_1} \mathbb{P} \nabla\cdot (\bw \otimes \bw - \bB \otimes \bB)\bigr\|_2 \, \mathrm{d}\tau,
\\
\|\bB(s)\|_2 &\leq \|e^{-sL_2} \bB_0\|_2 + \int_0^s \bigl\|e^{-(s-\tau)L_2} \mathbb{P} \nabla\cdot (\bw \otimes \bB - \bB \otimes \bw)\bigr\|_2 \, \mathrm{d}\tau.
\end{align*}
Starting from this consideration, the proofs of the main theorems rely on the strong energy inequality and on the following auxiliary proposition. Its first part provides an estimate for the linear terms, while the second part handles the nonlinear terms.

\begin{proposition} \label{lemma:a}
Let $L_i$, with $i =1,\,2$, be the linear operators defined in \cref{eqn:linearoperator}. Then the following propositions hold.
\begin{propenum}
\item \label{it:lma1} There exists a constant $\delta_3>0$ such that if $0<|\bb|\le \delta_3$, then for every $\Bf\in L^2_{\sigma}$,
\begin{align}\label{eqn:linearpart}
    \lim\limits_{t\to\infty} \frac{2}{t} \int_{\frac{t}{2}}^t \|e^{-sL_i}\Bf\|_2 \, \mathrm d s=0.
\end{align}    
 \item \label{it:lma2} Let $\frac{6}{5}<r\le2$ and let $\delta_2 = \delta_2(r) > 0$ be the constant from \cref{it:lmdecay2}. If $0<|\bb|\le \delta_2$, then for all $\bu,\bv\in H^1(\R^3)$,
    \begin{align}\label{eqn:nonlinearpart}
    \|e^{-sL_i} \mathbb{P} \nabla\cdot (\bu\otimes \bv)\|_2 \le C s^{-\frac{3}{2}\left(\frac{1}{r}-\frac{1}{2}\right)-\frac{1}{2}} 
    (\|\bu\|_2\|\bv\|_2)^{\frac{3}{2r}-\frac{1}{2}} (\|\nabla\bu\|_2\|\nabla \bv\|_2)^{\frac{3}{2}(1-\frac{1}{r})}.
    \end{align}
\end{propenum}
\end{proposition}

\begin{proof} The two claims are consequences of \cref{prop:decay}.
\begin{enumerate}[label={(\roman*)}]
    \item Let $\delta_3=\min\{\delta_1(\frac{3}{2}),\delta_1(2)\}$, where $\delta_1(r)$ is the constant defined in \cref{it:lmdecay1}. 
For all $\Bf \in L^2_\sigma$ and every $\varepsilon > 0$, we can choose $\Bf_\varepsilon \in C_{c,\sigma}^{\infty}$ such that $\|\Bf - \Bf_\varepsilon\|_2 \le \varepsilon$. Then 
\[
\|e^{-sL_i} \Bf\|_2 
\le \|e^{-sL_i} (\Bf - \Bf_\varepsilon)\|_2 + \|e^{-sL_i} \Bf_\varepsilon\|_2 
\le C \varepsilon + C s^{-\frac{1}{4}} \|\Bf_\varepsilon\|_{\frac32},
\]
where in the last inequality we used \cref{eqn:decay1}. Therefore we have
\[
\frac{2}{t} \int_{\frac{t}{2}}^{t} \|e^{-sL_i} \Bf\|_2 \, \mathrm d s 
\le C\varepsilon + Ct^{-\frac{1}{4}} \|\Bf_\varepsilon\|_{\frac{3}{2}}.
\]
From this, we deduce that
\[
\lim_{t \to \infty} \frac{2}{t} \int_{\frac{t}{2}}^{t} \|e^{-sL_i} \Bf\|_2 \, \mathrm ds \le C \varepsilon.
\]
Since $\varepsilon$ is arbitrary, it follows that $\lim\limits_{t \to \infty} \frac{2}{t} \int_{\frac{t}{2}}^{t} \|e^{-sL_i} \Bf\|_2 \, \mathrm ds = 0$.\\

\item  Using the decay property of $e^{-tL_i}$ in \cref{it:lmdecay2}, we have the estimate
\[
\|e^{-sL_i} \mathbb{P} \nabla\cdot (\bu \otimes \bv)\|_2 
\leq C s^{-\frac{3}{2}\left(\frac{1}{r}-\frac{1}{2}\right) - \frac{1}{2}} \|\bu \otimes \bv\|_r.
\]
H\"older's inequality and Gagliardo--Nirenberg's inequality imply 
\[
\|\bu \otimes \bv\|_r 
\le \|\bu\|_{2r} \| \bv\|_{2r}
\le (\|\bu\|_{2}\|\bv\|_{2})^{1-\theta} 
(\|\nabla \bu\|_{2}\|\nabla \bv\|_{2})^{\theta},
\]
where $\theta$ satisfies $\frac{1}{2r}=\frac{1-\theta}{2}+\theta(\frac{1}{2}-\frac{1}{3})$; in particular, $\theta\in(0,1)$ precisely when $r\in(\frac{6}{5},2]$. Consequently,
\[
\|\bu \otimes \bv\|_r 
\le  (\|\bu\|_{2}\|\bv\|_{2})^{\frac{3}{2r}-\frac{1}{2}} 
(\|\nabla \bu\|_{2}\|\nabla \bv\|_{2})^{\frac{3}{2}(1-\frac{1}{r})},
\]
which completes the proof.
\end{enumerate}
\end{proof}

We now turn to the proof of \cref{thm:stability}.  

\begin{proof}[Proof of \cref{thm:stability}] 

\uline{Step 1.} \emph{Duhamel-type representation.} 
%Let $0<|\bb|\le \min\{\ve_0,\frac{1}{6C_0} \}$. 
As $(\bw,\bB)$ is a weak solution to \cref{eqn:PerturbedEqn}, we have
\begin{align}\label{eqn:weakformulation1}
\langle \bw(s), \boldsymbol{\varphi}(s)  \rangle
+\int_0^s \Big[ 
\langle \bw\cdot  \nabla \bw, \boldsymbol{\varphi} \rangle
- \langle \bB\cdot  \nabla \bB, \boldsymbol{\varphi} \rangle
+ a_1(\bw,\boldsymbol{\vp})
\Big]\, \mathrm d\tau 
= \langle \bw(0), \boldsymbol{\varphi}(0)  \rangle 
+ \int_0^s \langle \bw, \boldsymbol{\varphi}_\tau \rangle \, \mathrm d\tau
\end{align}
for all $s\geq 0$ and all $\boldsymbol{\varphi}$ satisfying 
\[
\boldsymbol{\varphi} \in C\big( [0, \infty); L^2_\sigma \cap L^3 \big), \quad 
\nabla \boldsymbol{\varphi} \in L^2_{\mathrm{loc}}\big( [0,\infty); L^2 \big), \quad 
\partial_\tau \boldsymbol{\varphi} \in L^2_{\mathrm{loc}}\big( [0,\infty); (H^1_\sigma)^* \big).
\]


% Taking $\boldsymbol{\vp}(\tau)=e^{-(s-\tau)L_1^*}\boldsymbol{\psi}$ as a test function, where $\boldsymbol{\psi}\in C^{\infty}_{c,\sigma}$, we get
% \begin{align*}
% \langle \bw(s), \psi  \rangle
% = \langle e^{-sL_1}\bw_0,  \psi  \rangle 
% - \int_0^s  
% \langle e^{-(s-\tau)L_1}\mathbb{P}( \bw\cdot \nabla \bw- \bB\cdot  \nabla \bB), \psi \rangle \, \mathrm d\tau.
% \end{align*}
% Here we used 
% \[
% \langle \bw, \boldsymbol{\varphi}_\tau \rangle=\langle \bw, L_1^*\boldsymbol{\varphi}  \rangle =a_1(\bw, \boldsymbol{\varphi}).
% \]
% For sufficiently small $|\bb|$, we consider the test function $\boldsymbol{\vp}(\tau)=e^{-(s-\tau)L_1^*}\boldsymbol{\psi}$ with $\boldsymbol{\psi}\in C^{\infty}_{c,\sigma}$.  
% We first verify that $\boldsymbol{\vp}$ belongs to the admissible class of test functions defined in \eqref{eqn:testfcnspace}.  
There exists a constant $0<\delta_4\le \min \{ \ve_0, \frac{1}{6C_0}\}$ such that for any $|\bb|\le\delta_4$ and $\boldsymbol{\psi}\in C^{\infty}_{c,\sigma}$, the function $\boldsymbol{\vp}(\tau)=e^{-(s-\tau)L_i^*}\boldsymbol{\psi}$ belongs to the admissible test function space for $i=1,2$.
Indeed, by the analyticity of the semigroup $e^{-tL_1^*}$ (see, \cref{thm:analyticsemigroup} and \cite[Theorem 3.1]{LZZ23}), we have $\boldsymbol{\vp}\in C([0,\infty);L^2_\sigma\cap L^3)$.  
Moreover, using the commutativity $L_1^{*\frac12}e^{-tL_1^*}=e^{-tL_1^*}L_1^{*\frac12}$ (see, e.\,g., \cite[Proposition 3.1.1(f)]{Haase06} and \cite[Theorem 1.2.4(c)]{Pazy1983}) together with the estimates \cref{eqn:corollary1} and \cref{eqn:corollary2}, we obtain
\[
\| \nabla \boldsymbol{\vp}(\tau) \|_2
\lesssim \| L_1^{*\frac{1}{2}} e^{-(s-\tau)L_1^*} \boldsymbol{\psi} \|_2  
= \| e^{-(s-\tau)L_1^*} L_1^{*\frac{1}{2}} \boldsymbol{\psi} \|_2  
\le \| L_1^{*\frac{1}{2}} \boldsymbol{\psi} \|_2 
\lesssim \|\nabla \boldsymbol{\psi} \|_2,
\]
which implies $\nabla\boldsymbol{\vp}\in L^\infty_{\mathrm{loc}}([0,\infty);L^2)\subset L^2_{\mathrm{loc}}([0,\infty);L^2)$.  
For the time derivative, we note that $\partial_\tau\boldsymbol{\vp}(\tau)=L_1^*\boldsymbol{\vp}(\tau)$. Then for any $\bv\in H^1_\sigma$,
\[
|\langle\partial_\tau\boldsymbol{\vp}(\tau),\bv\rangle|=|a_1(\boldsymbol{\vp}(\tau),\bv)|\le C\|\boldsymbol{\vp}(\tau)\|_{H^1}\|\bv\|_{H^1},
\]
so that by duality $\|\partial_\tau\boldsymbol{\vp}(\tau)\|_{(H^1_\sigma)^*}\le C(\|\boldsymbol{\vp}(\tau)\|_2+\|\nabla\boldsymbol{\vp}(\tau)\|_2)$.  
The boundedness of $\|\boldsymbol{\vp}(\tau)\|_2$ (from \cref{eqn:corollary2}) and of $\|\nabla\boldsymbol{\vp}(\tau)\|_2$ then yields $\partial_\tau\boldsymbol{\vp}\in L^2_{\mathrm{loc}}([0,\infty);(H^1_\sigma)^*)$.  
Thus $\boldsymbol{\vp}$ satisfies the required regularity conditions.

Taking this $\boldsymbol{\vp}$ as a test function in \eqref{eqn:weakformulation1} and using $\langle\bw,\partial_\tau\boldsymbol{\vp}\rangle=\langle\bw,L_1^*\boldsymbol{\vp}\rangle=a_1(\bw,\boldsymbol{\vp})$, we obtain
\begin{align*}
\langle \bw(s), \boldsymbol{\psi} \rangle
= \Big\langle e^{-sL_1}\bw_0
- \int_0^s e^{-(s-\tau)L_1}\mathbb{P}(\bw\cdot\nabla\bw-\bB\cdot\nabla\bB)\, \mathrm d\tau, \boldsymbol{\psi} \Big\rangle.
\end{align*}
By duality, this gives
\begin{align}\label{eqn:estimate2-1}
\|\bw(s)\|_2 \leq \|e^{-sL_1} \bw_0\|_2 + \int_0^s \bigl\|e^{-(s-\tau)L_1} \mathbb{P} \nabla\cdot (\bw \otimes \bw - \bB \otimes \bB)\bigr\|_2 \, \mathrm{d}\tau.
\end{align}
Similarly, taking $\boldsymbol{\vp}(\tau)=e^{-(s-\tau)L_2^*}\boldsymbol{\psi}$ as a test function in \cref{eqn:weak}$_2$ yields
\begin{align}\label{eqn:estimate2-2}
\|\bB(s)\|_2 \leq \|e^{-sL_2} \bB_0\|_2 + \int_0^s \bigl\|e^{-(s-\tau)L_2} \mathbb{P} \nabla\cdot (\bw \otimes \bB - \bB \otimes \bw)\bigr\|_2 \, \mathrm{d}\tau.
\end{align}



\uline{Step 2.} \emph{Energy monotonicity and time averaging}. If $0<|\bb|<\min\{\ve_0,\frac{1}{4C_0}\}$, Hardy's inequality \cref{eqn:Hardy} and the strong energy inequality \cref{eqn:energyineq1} imply
\begin{align}\label{eqn:energyineq2}
\|\bw(t)\|_2^2 + \|\bB(t)\|_2^2 + 2(1-4C_0|\bb|) \int_s^t \left(\|\nabla \bw\|_2^2 + \|\nabla \bB\|_2^2\right) \, \mathrm d\tau \leq \|\bw(s)\|_2^2 + \|\bB(s)\|_2^2,
\end{align}
which yields
\begin{align}\label{eq:s-est-0}
\|\bw(t)\|_2^2 + \|\bB(t)\|_2^2 \leq \|\bw(s)\|_2^2 + \|\bB(s)\|_2^2
\end{align}
for almost every $s\geq 0$ (including $s=0$) and for all $t\geq s$.
Therefore, we obtain
\begin{align*}
\frac{\|\bw(t)\|_2 + \|\bB(t)\|_2}{\sqrt{2}} 
\leq \sqrt{\|\bw(t)\|_2^2 + \|\bB(t)\|_2^2} 
\leq \sqrt{\|\bw(s)\|_2^2 + \|\bB(s)\|_2^2} 
\leq \|\bw(s)\|_2 + \|\bB(s)\|_2.
\end{align*}
%In particular, we note that 
%\begin{align}\label{eq:s=0}
% \frac{\|\bw(t)\|_2+\|\bB(t)\|_2}{\sqrt{2}} \le \|\bw_0\|_2 +\|\bB_0\|_2 .  
%\end{align}
Integrating the inequality over $s\in[\frac{t}{2}, t]$ gives
\begin{align}\label{eqn:estimate1}
\frac{\|\bw(t)\|_2 + \|\bB(t)\|_2}{\sqrt{2}} 
\leq \frac{2}{t} \int_{\frac{t}{2}}^{t} \|\bw(s)\|_2 \, \mathrm ds + \frac{2}{t} \int_{\frac{t}{2}}^{t} \|\bB(s)\|_2 \, \mathrm ds. 
\end{align}

\uline{Step 3.} \emph{Combination of linear and nonlinear contributions.} Combining \cref{eqn:estimate2-1}, \cref{eqn:estimate2-2} and \cref{eqn:estimate1}, we deduce 
\begin{equation}\label{eqn:estimate3}
\begin{aligned}
\frac{\|\bw(t)\|_2 + \|\bB(t)\|_2}{\sqrt{2}} 
 &\leq \frac{2}{t} \int_{\frac{t}{2}}^t \bigl( \|e^{-sL_1} \bw_{0}\|_2 + \|e^{-sL_2} \bB_0\|_2 \bigr) \, \mathrm{d} s\\
& +  \frac{2}{t} \int_{\frac{t}{2}}^{t} \Bigg[ \int_0^s \Bigl( \|e^{-(s-\tau)L_1} \mathbb{P}\nabla\cdot (\bw \otimes \bw - \bB \otimes \bB)\|_2 
\\ & \qquad \qquad \qquad  + \|e^{-(s-\tau)L_2} \mathbb{P}\nabla\cdot (\bw\otimes \bB - \bB \otimes \bw)\|_2 \Bigr) \, \mathrm{d} \tau \Bigg] \mathrm{d}s.
\end{aligned}
\end{equation}

\uline{Step 4.} \emph{Linear contribution.} If $|\bb|\le \delta_3$, where $\delta_3$ is the constant in \cref{it:lma1} ,
\[
\lim\limits_{t\to \infty} \frac{2}{t} \int_{\frac{t}{2}}^t \bigl( \|e^{-sL_1} \bw_{0}\|_2 + \|e^{-sL_2} \bB_0\|_2 \bigr) \, \mathrm{d} s=0.
\]

\uline{Step 5.} \emph{Nonlinear contribution.} Inequality \cref{eqn:nonlinearpart} implies that
\begin{align*}
&\frac{2}{t} \int_{\frac{t}{2}}^{t} \left[ \int_0^s  \|e^{-(s-\tau)L_1} \mathbb{P}\nabla\cdot (\bw \otimes \bw)\|_2 \, \mathrm{d} \tau \right] \mathrm{d}s\\
& \le  \frac{C}{t} \int_{\frac{t}{2}}^{t} 
\left[ \int_0^s (s-\tau)^{-\frac{3}{2}(\frac{1}{r_1}-\frac{1}{2})-\frac{1}{2}} \|\bw(\tau)\|_2^{\frac{3}{r_1}-1}  h^{\frac
32(1-\frac{1}{r_1})} \, \mathrm{d} \tau \right] \mathrm{d}s,
%& \le   \frac{C}{t} \int_{\frac{t}{2}}^{t} \left[ \int_0^s (s-\tau)^{-\frac{3}{2}(\frac{1}{r_1}-\frac{1}{2})-\frac{1}{2}}  \|\nabla \bw(\tau)\|_2^{3(1-\frac{1}{r_1})} \, \mathrm{d} \tau \right] \mathrm{d}s,
\end{align*}
where $h(\tau)=\|\nabla \bw(\tau)\|_2^2$ and $r_1\in (\frac{6}{5},2]$ is a constant to be determined. 
\cref{eqn:energyineq2} implies $h\in L^{1}(0,\infty)$; moreover, by \eqref{eq:s-est-0}, $\|\bw(\tau)\|_2 \le \|\bw_0\|_2+\|\bB_0\|_2$.
Therefore
\[
\frac{C}{t} \int_{\frac{t}{2}}^{t} 
\left[ \int_0^s (s-\tau)^{-\frac{3}{2}(\frac{1}{r_1}-\frac{1}{2})-\frac{1}{2}} \|\bw(\tau)\|_2^{\frac{3}{r_1}-1}  h^{\frac
32(1-\frac{1}{r_1})} \, \mathrm{d} \tau \right] \mathrm{d}s
\le 
\frac{C}{t} \int_{\frac{t}{2}}^{t} 
|\cdot|^{-\frac{3}{2}(\frac{1}{r_1}-\frac{1}{2})-\frac{1}{2}} \ast h^{\frac{3}{2}(1-\frac{1}{r_1})}\, \mathrm{d}s.
\]
Using \cref{eqn:weakLpnorm}, we have
\begin{align*}
\frac{C}{t} \int_{\frac{t}{2}}^{t} 
|\cdot|^{-\frac{3}{2}(\frac{1}{r}-\frac{1}{2})-\frac{1}{2}} \ast h^{\frac{3}{2}(1-\frac{1}{r_1})}\, \mathrm{d}s
& \le Ct^{-\frac{1}{\alpha_1}}\left\| |\cdot|^{-\frac{3}{2}(\frac{1}{r_1}-\frac{1}{2})-\frac{1}{2}} \ast h^{\frac{3}{2}(1-\frac{1}{r_1})} \right\|_{\alpha_1,\infty}\\
& \le  Ct^{-\frac{1}{\alpha_1}}
\left\| |\cdot|^{-\frac{3}{2}(\frac{1}{r_1}-\frac{1}{2})-\frac{1}{2}}\right\|_{c_1,\infty} 
\left\| h^{\frac{3}{2}(1-\frac{1}{r_1})} \right\|_{c_2}.
\end{align*}
Here we used \cref{eqn:Youngineq} and the embedding $L^p \subset L^{p,\infty}$ in the second line, and $\alpha_1,c_1,c_2>1$ are constants satisfying
\[
\frac{1}{c_1}=\frac{3}{2}\left(\frac{1}{r_1}-\frac{1}{2}\right)+\frac{1}{2},\quad
\frac{1}{c_2}=\frac{3}{2}\left(1-\frac{1}{r_1}\right),
\]
and
\[
1+\frac{1}{\alpha_1}=\frac{1}{c_1}+\frac{1}{c_2}.
\]
Note that $\frac{1}{c_1}+\frac{1}{c_2}=\frac{5}{4}$, so $\frac{1}{\alpha_1}=\frac{1}{4}$ provided $c_1,c_2\in (1,\infty)$. For instance, taking $r_1=\frac{3}{2}$ gives $c_1,c_2\in(1,\infty)$.
Therefore, if $|\bb|\le \delta_2(\frac{3}{2})$ with $\delta_2(\frac{3}{2})$ as in \cref{it:lma2}, 
\[
\frac{2}{t} \int_{\frac{t}{2}}^{t} \left[ \int_0^s \|e^{-(s-\tau)L_1} \mathbb{P}\nabla\cdot (\bw \otimes \bw)\|_2 \, \mathrm{d} \tau \right] \mathrm{d}s \le C t^{-\frac{1}{4}}.
\]
Similarly, 
\begin{align*}
&\frac{2}{t} \int_{\frac{t}{2}}^{t} \left[ \int_0^s \Bigl( \|e^{-(s-\tau)L_1} \mathbb{P}\nabla\cdot (\bw \otimes \bw - \bB \otimes \bB)\|_2 
+ \|e^{-(s-\tau)L_2} \mathbb{P}\nabla\cdot (\bw\otimes \bB - \bB \otimes \bw)\|_2 \Bigr) \, \mathrm{d} \tau \right] \, \mathrm{d}s \\
&\le Ct^{-\frac{1}{4}}
\end{align*}
provided $|\bb| \le \delta_2(\frac{3}{2})$.


% where we used $\|\bw(\tau)\|_2 \le \|\bw_0\|_2+\|\bB_0\|_2$ according to the strong energy inequality \eqref{eq:s-est-0}, and $r_1\in (\frac{6}{5},2]$ is a constant to be determined. \cref{eqn:energyineq2} implies $\|\nabla \bw\|_2^2\in L^{1}(0,\infty)$. Denote by $h(\tau)=\|\nabla \bw(\tau)\|_2^2$. Then
% \[
% \frac{C}{t} \int_{\frac{t}{2}}^{t} 
% \left[ \int_0^s (s-\tau)^{-\frac{3}{2}(\frac{1}{r_1}-\frac{1}{2})-\frac{1}{2}}  \|\nabla \bw(\tau)\|_2^{3(1-\frac{1}{r_1})} \, \mathrm{d} \tau \right] \mathrm{d}s 
% \le 
% \frac{C}{t} \int_{\frac{t}{2}}^{t} 
% |\cdot|^{-\frac{3}{2}(\frac{1}{r_1}-\frac{1}{2})-\frac{1}{2}} \ast h^{\frac{3}{2}(1-\frac{1}{r_1})}\, \mathrm{d}s.
% \]
% Using \cref{eqn:weakLpnorm}, we have
% \begin{align*}
% \frac{C}{t} \int_{\frac{t}{2}}^{t} 
% |\cdot|^{-\frac{3}{2}(\frac{1}{r}-\frac{1}{2})-\frac{1}{2}} \ast h^{\frac{3}{2}(1-\frac{1}{r_1})}\, \mathrm{d}s
% & \le Ct^{-\frac{1}{\alpha_1}}\left\| |\cdot|^{-\frac{3}{2}(\frac{1}{r_1}-\frac{1}{2})-\frac{1}{2}} \ast h^{\frac{3}{2}(1-\frac{1}{r_1})} \right\|_{\alpha_1,\infty}\\
% & \le  Ct^{-\frac{1}{\alpha_1}}
% \left\| |\cdot|^{-\frac{3}{2}(\frac{1}{r_1}-\frac{1}{2})-\frac{1}{2}}\right\|_{c_1,\infty} 
% \left\| h^{\frac{3}{2}(1-\frac{1}{r_1})} \right\|_{c_2}.
% \end{align*}
% Here we used Young's inequality \cref{eqn:Youngineq} and the embedding $L^p \subset L^{p,\infty}$ in the second line, and $\alpha_1,c_1,c_2>1$ are constants satisfying
% \[
% \frac{1}{c_1}=\frac{3}{2}\left(\frac{1}{r_1}-\frac{1}{2}\right)+\frac{1}{2},\quad
% \frac{1}{c_2}=\frac{3}{2}\left(1-\frac{1}{r_1}\right),
% \]
% and
% \[
% 1+\frac{1}{\alpha_1}=\frac{1}{c_1}+\frac{1}{c_2}.
% \]
% Letting %$\frac{1}{c_1}=\frac{3}{2}(\frac{1}{r_1}-\frac{1}{2})+\frac{1}{2}$ and $\frac{1}{c_2}=\frac{3}{2}(1-\frac{1}{r_1})$ with 
% $r_1=\frac{3}{2}$, we have $c_1,c_2 \in (1,\infty)$ and $\frac{1}{\alpha_1}=\frac{1}{4}$. Therefore if $|\bb|\le \delta_2(\frac{3}{2})$, where $\delta_2(\frac{3}{2})$ is the constant in \cref{eqn:nonlinearpart}, 
% \[
% \frac{2}{t} \int_{\frac{t}{2}}^{t} \left[ \int_0^s \Bigl( \|e^{-(s-\tau)L_1} \mathbb{P}\nabla\cdot (\bw \otimes \bw)\|_2 \, \mathrm{d} \tau \right] \mathrm{d}s \le C t^{-\frac{1}{4}}.
% \]
% Similarly, 
% \begin{align*}
% &\frac{2}{t} \int_{\frac{t}{2}}^{t}  \int_0^s \Bigl( \|e^{-(s-\tau)L_1} \mathbb{P}\nabla\cdot (\bw \otimes \bw - \bB \otimes \bB)\|_2 
% + \|e^{-(s-\tau)L_2} \mathbb{P}\nabla\cdot (\bw\otimes \bB - \bB \otimes \bw)\|_2 \Bigr) \, \mathrm{d} \tau \, \mathrm{d}s \\
% &\le Ct^{-\frac{1}{4}}
% \end{align*}
% provided $|\bb| \le \delta_2(\frac{3}{2})$.

\uline{Step 6.} \emph{Conclusion.} In conclusion, if $0<|\bb|\le \delta$, where $\delta=\min\{ \delta_4, \delta_3, \delta_2(\frac{3}{2})\}$, \cref{eqn:estimate3} yields
\[
\lim\limits_{t\to \infty}\|\bw(t)\|_2 
=\lim\limits_{t\to \infty} \|\bB(t)\|_2=0.
\]
\end{proof}


We now establish the decay estimate for the perturbed equation \cref{eqn:PerturbedEqn}. The proof proceeds by iteration.

\begin{proof}[Proof of \cref{thm:decay}]

\uline{Step 1.} \emph{Linear decay estimate.} By
\cref{it:lmdecay1}
\[
\frac{2}{t} \int_0^t \bigl( \|e^{-sL_1} \bw_{0}\|_2 + \|e^{-sL_2} \bB_0\|_2 \bigr) \, \mathrm{d} s 
\le Ct^{-\frac{3}{2}(\frac{1}{q}-\frac{1}{2})} (\|\bw_0\|_q + \|\bB_0\|_q).
\]


\uline{Step 2.} \emph{Baseline decay bound.} 
From the estimates obtained in the proof of \cref{thm:stability}, 
\begin{align*}
&\frac{2}{t} \int_{\frac{t}{2}}^{t} \left[ \int_0^s \Bigl( \|e^{-(s-\tau)L_1} \mathbb{P}\nabla\cdot (\bw \otimes \bw - \bB \otimes \bB)\|_2 
+ \|e^{-(s-\tau)L_2} \mathbb{P}\nabla\cdot (\bw\otimes \bB - \bB \otimes \bw)\|_2 \Bigr) \, \mathrm{d} \tau \right] \, \mathrm{d}s \\
&\le Ct^{-\frac{1}{4}}.
\end{align*}
Combining these estimates with \cref{eqn:estimate3}, we have
\[
\|\bw(t)\|_2+\|\bB(t)\|_2 \le Ct^{-\frac{3}{2}(\frac{1}{q}-\frac{1}{2})} + Ct^{-\frac{1}{4}}.
\]

\uline{Step 3.} \emph{Decay regimes according to the integrability exponent $q$.} Thus \cref{thm:decay} holds whenever $-\frac{1}{4} \le -\frac{3}{2}(\frac{1}{q}-\frac{1}{2})$, that is, $\frac{3}{2} \le q <2$. If $1 \le q < \frac{3}{2}$, $\|\bw(t)\|_2+\|\bB(t)\|_2 \le Ct^{-\frac{1}{4}}$. 


\uline{Step 4.} \emph{Bootstrap improvement of decay for $1 \le q < \frac{3}{2}$.}
Similarly, \cref{eqn:nonlinearpart} shows
\begin{align*}
&\frac{2}{t} \int_{\frac{t}{2}}^{t} \left[ \int_0^s  \|e^{-(s-\tau)L_1} \mathbb{P}\nabla\cdot (\bw \otimes \bw)\|_2 \, \mathrm{d} \tau \right] \mathrm{d}s\\
& \le  \frac{C}{t} \int_{\frac{t}{2}}^{t} 
\left[ \int_0^s (s-\tau)^{-\frac{3}{2}(\frac{1}{r_2}-\frac{1}{2})-\frac{1}{2}} \|\bw(\tau)\|_2^{\frac{3}{r_2}-1}  h^{\frac32(1-\frac{1}{r_2})} \, \mathrm{d} \tau \right] \mathrm{d}s.
\end{align*}
Here $h(\tau)=\|\bw(\tau)\|_2^2$, and $r_2\in (\frac{6}{5},2]$ is a constant to be determined. Using $\|\bw(\tau)\|_2 \le C\tau^{-\frac{1}{4}}$, we obtain
\begin{align*}
&\frac{C}{t} \int_{\frac{t}{2}}^{t} 
\left[ \int_0^s (s-\tau)^{-\frac{3}{2}(\frac{1}{r_2}-\frac{1}{2})-\frac{1}{2}} \|\bw(\tau)\|_2^{\frac{3}{r_2}-1}  h^{\frac32(1-\frac{1}{r_2})} \, \mathrm{d} \tau \right] \mathrm{d}s\\
& \le   \frac{C}{t} \int_{\frac{t}{2}}^{t} 
\left[ \int_0^s (s-\tau)^{-\frac{3}{2}(\frac{1}{r_2}-\frac{1}{2})-\frac{1}{2}}  \tau^{-\frac{1}{4}(\frac{3}{r_2}-1)}  h^{\frac{3}{2}(1-\frac{1}{r_2})} \, \mathrm{d} \tau \right] \mathrm{d}s\\
&=\frac{C}{t} \int_{\frac{t}{2}}^{t} 
|\cdot|^{-\frac{3}{2}(\frac{1}{r_2}-\frac{1}{2})-\frac{1}{2}} \ast ( |\cdot|^{-\frac{1}{4}(\frac{3}{r_2}-1)} \, h^{\frac{3}{2}(1-\frac{1}{r_2})})\, \mathrm{d}s,
\end{align*}
% Using $\|\bw(\tau)\|_2 \le C\tau^{-\frac{1}{4}}$, we have
% \begin{align*}
% &\frac{2}{t} \int_{\frac{t}{2}}^{t} \left[ \int_0^s \Bigl( \|e^{-(s-\tau)L_1} \mathbb{P}\nabla\cdot (\bw \otimes \bw)\|_2 \, \mathrm{d} \tau \right] \mathrm{d}s\\
% & \le  \frac{C}{t} \int_{\frac{t}{2}}^{t} 
% \left[ \int_0^s (s-\tau)^{-\frac{3}{2}(\frac{1}{r_2}-\frac{1}{2})-\frac{1}{2}} \|\bw(\tau)\|_2^{\frac{3}{r_2}-1}  \|\nabla \bw(\tau)\|_2^{3(1-\frac{1}{r_2})} \, \mathrm{d} \tau \right] \mathrm{d}s\\
% & \le   \frac{C}{t} \int_{\frac{t}{2}}^{t} 
% \left[ \int_0^s (s-\tau)^{-\frac{3}{2}(\frac{1}{r_2}-\frac{1}{2})-\frac{1}{2}}  \tau^{-\frac{1}{4}(\frac{3}{r_2}-1)}  h^{\frac{3}{2}(1-\frac{1}{r_2})} \, \mathrm{d} \tau \right] \mathrm{d}s\\
% &=\frac{C}{t} \int_{\frac{t}{2}}^{t} 
% |\cdot|^{-\frac{3}{2}(\frac{1}{r_2}-\frac{1}{2})-\frac{1}{2}} \ast ( |\cdot|^{-\frac{1}{4}(\frac{3}{r_2}-1)} \, h^{\frac{3}{2}(1-\frac{1}{r_2})})\, \mathrm{d}s,
% \end{align*}
% Here $r_2\in (\frac{6}{5},2]$ is a constant to be determined, and $h(\tau)=\|\bw(\tau)\|_2^2$. 
Recalling
\cref{eqn:weakLpnorm}, \cref{eqn:Youngineq}, and \cref{eqn:Holderineq}, we deduce 
\begin{align*}
&\frac{C}{t} \int_{\frac{t}{2}}^{t} 
|\cdot|^{-\frac{3}{2}(\frac{1}{r_2}-\frac{1}{2})-\frac{1}{2}} \ast ( |\cdot|^{-\frac{1}{4}(\frac{3}{r_2}-1)} \, h^{\frac{3}{2}(1-\frac{1}{r_2})})\, \mathrm{d}s \\
& \le 
C t^{-\frac{1}{\alpha_2}}
\left\| |\cdot|^{-\frac{3}{2}(\frac{1}{r_2}-\frac{1}{2})-\frac{1}{2}} \ast ( |\cdot|^{-\frac{1}{4}(\frac{3}{r_2}-1)} \, h^{\frac{3}{2}(1-\frac{1}{r_2})}) \right\|_{\alpha_2,\infty}\\
& \le  C t^{-\frac{1}{\alpha_2}}
\left\| |\cdot|^{-\frac{3}{2}(\frac{1}{r_2}-\frac{1}{2})-\frac{1}{2}} \right\|_{c_3,\infty}
\left\| |\cdot|^{-\frac{1}{4}(\frac{3}{r_2}-1)} \right\|_{c_4,\infty}
\left\| h^{\frac{3}{2}(1-\frac{1}{r_2})} \right\|_{c_5,\infty},
\end{align*}
where $\alpha_2,c_3,c_4,c_5>1$ are constants satisfying
\[
\frac{1}{c_3}=\frac{3}{2}\left(\frac{1}{r_2}-\frac{1}{2}\right)+\frac{1}{2},\quad
\frac{1}{c_4}= \frac{1}{4}\left(\frac{3}{r_2}-1\right),\quad 
\frac{1}{c_5}=\frac{3}{2}\left(1-\frac{1}{r_2}\right),
\]
and
\[
0< \frac{1}{c_4}+\frac{1}{c_5} <1,\quad 1+\frac{1}{\alpha_2}=\frac{1}{c_3}+\frac{1}{c_4}+\frac{1}{c_5}.
\]
Taking $r_2=\frac{15}{8}$ gives $\frac{1}{\alpha_2}=\frac{2}{5}$, which yields
\[
\|\bw(t)\|_2+\|\bB(t)\|_2 \le Ct^{-\frac{3}{2}(\frac{1}{q}-\frac{1}{2})} + Ct^{-\frac{2}{5}}.
\]
Thus \cref{thm:decay} holds provided $-\frac{2}{5} \le -\frac{3}{2}(\frac{1}{q}-\frac{1}{2})$, i.\,e., $\frac{30}{23} \le q <\frac{3}{2}$.

\uline{Step 
5.} \emph{Final bootstrap and completion of the proof.} If $1 \le q < \frac{30}{23}$, $\|\bw(t)\|_2+\|\bB(t)\|_2 \le Ct^{-\frac{2}{5}}$. Analogously,
\begin{align*}
&\frac{2}{t} \int_{\frac{t}{2}}^{t} \left[ \int_0^s \|e^{-(s-\tau)L_1} \mathbb{P}\nabla\cdot (\bw \otimes \bw)\|_2 \, \mathrm{d} \tau \right] \mathrm{d}s\\
& \le  C t^{-\frac{1}{\alpha_3}}
\left\| |\cdot|^{-\frac{3}{2}(\frac{1}{r_3}-\frac{1}{2})-\frac{1}{2}} \right\|_{c_6,\infty}
\left\| |\cdot|^{-\frac{2}{5}(\frac{3}{r_3}-1)} \right\|_{c_7,\infty}
\left\| h^{\frac{3}{2}(1-\frac{1}{r_3})} \right\|_{c_8,\infty}.
\end{align*}
Here $r_3\in (\frac{6}{5},2]$ and $\alpha_3,c_6,c_7,c_8>1$ are constants satisfying
\[
\frac{1}{c_6}=\frac{3}{2}\left(\frac{1}{r_3}-\frac{1}{2}\right)+\frac{1}{2},\quad
\frac{1}{c_7}= \frac{2}{5}\left(\frac{3}{r_3}-1\right),\quad 
\frac{1}{c_8}=\frac{3}{2}\left(1-\frac{1}{r_3}\right),
\]
and
\[
0< \frac{1}{c_7}+\frac{1}{c_8} <1, \quad 1+\frac{1}{\alpha_3}=\frac{1}{c_6}+\frac{1}{c_7}+\frac{1}{c_8}.
\]
Choosing $r_3=\frac{4}{3}$ gives $\frac{1}{\alpha_3}=\frac{3}{4}$, which yields
\[
\|\bw(t)\|_2+\|\bB(t)\|_2 \le Ct^{-\frac{3}{2}\left(\frac{1}{q}-\frac{1}{2}\right)} + Ct^{-\frac{3}{4}}.
\]
Thus the proof is complete.
\end{proof}

\begin{remark}[Refinement of the decay rate]
In the proof of \cref{thm:decay}, the constants $r_2$ and $r_3$ can be chosen so that $r_2 > 3/2$ is arbitrarily close to $3/2$ and $r_3 > \frac{6}{5}$ is arbitrarily close to $\frac{6}{5}$. With such a choice, we have $\frac{1}{\alpha_3} = 1 - \varepsilon$ for an arbitrarily small $\varepsilon > 0$, which gives the decay estimate
\[
\|\bw(t)\|_2+\|\bB(t)\|_2 \le C t^{-\frac{3}{2}\left(\frac{1}{q}-\frac{1}{2}\right)} + C t^{-(1-\varepsilon)}.
\]
\end{remark}

\vspace{0.5cm}

\section*{Acknowledgments}

N.~De Nitti is a member of the Gruppo Nazionale per l'Analisi Matematica, la Probabilità e le loro Applicazioni (GNAMPA) of the Istituto Nazionale di Alta Matematica (INdAM) and has received support from the INdAM--GNAMPA Project 2026 \textit{Modelli Non-locali in Fluidodinamica, Traffico ed Elasticità} (CUP:~E53C25002010001). He is grateful to M.~Sammartino and V.~Sciacca for helpful discussions on MHD models at the University of Palermo. He acknowledges Soochow University, Duke Kunshan University, and Shanghai Jiao Tong University for their kind hospitality.

Y.~Wang is partially supported by NSFC grant 12271389, and by the Natural Science Foundation of Jiangsu Province (grant~BK20240147).

S.~Zhang is partially supported by the Postgraduate Research \&
Practice Innovation Program of Jiangsu Province (grant~KYCX24\_3285). He acknowledges the kind hospitality of Duke Kunshan University and Shanghai Jiao Tong University.




\vspace{0.5cm}

\printbibliography

\vfill 

\end{document}
